\documentclass[UTF8,a4paper,12pt]{article}
\usepackage{ctex}
\usepackage{geometry}
\usepackage{enumitem}
\usepackage{hyperref}
\usepackage{xcolor} % 引入颜色包实现红色文字
\geometry{left=2.5cm,right=2.5cm,top=2.5cm,bottom=2.5cm}
\setlist[enumerate]{itemsep=0.5ex,parsep=0.3ex}
\setlist[itemize]{itemsep=0.5ex,parsep=0.3ex}
% 封面按要求修改
\title{\bfseries 《走进新能源》期中报告}
\author{}
\date{2026年4月}
\begin{document}
\maketitle

% 醒目红色加粗提示
\noindent{\Huge\bfseries\color{red}重要声明：本文仅作课程写作参考，严禁直接抄袭！}\\[2em]

\begin{abstract}
作为电气工程相关专业学生，通过《走进新能源》课程学习，我对全球能源转型有了系统认识。本文围绕课程作业要求，结合课堂所学，用通俗易懂的语言，梳理了化石能源、核能、可再生能源、可燃冰、新型储能、新能源前沿及能源政策等核心内容，分析了新能源发展的现状、挑战与未来方向，形成本份期中报告。
\end{abstract}

\section{引言}
能源是日常生活与工业生产的基础，过去百年，石油、天然气、煤炭等化石能源支撑了全球经济发展，但也带来了全球变暖、环境污染、能源安全等问题。随着我国“双碳”目标的提出，全球能源正加速向清洁低碳转型。作为电气专业学生，深入理解新能源发展逻辑，是未来参与新型电力系统建设的基础，本报告结合课程知识，对新能源相关核心议题展开分析。

\section{化石能源的未来：从主导到转型}
\subsection{二氧化碳的影响与“石油峰值”}
化石能源燃烧产生的二氧化碳是温室效应的核心来源，导致全球气温升高、极端天气频发，威胁生态与人类生存。“石油峰值”理论指出，石油作为不可再生资源，产量终将达到峰值后持续下降，叠加地缘政治冲突，能源安全问题日益突出，倒逼全球加速向新能源转型。

\subsection{化石能源的转型方向}
化石能源不会立即淘汰，而是向清洁化、高效化转型：一是通过超超临界技术提升煤电效率，减少排放；二是通过碳捕集（CCUS）封存二氧化碳，实现近零排放；三是作为新能源的调峰电源，保障电网稳定；四是从燃料向化工原料转型，提升资源价值。

\section{核能：安全清洁的基荷能源}
\subsection{核能的定位与快堆技术}
核能是唯一可大规模替代化石能源的基荷清洁能源，能量密度高、零碳排放、供电稳定。短期可替代煤电助力碳达峰，长期是碳中和的核心支撑。快中子反应堆能将铀-238转化为钚-239，将天然铀利用率从1\%提升至60\%以上，解决核燃料短缺问题，我国霞浦快堆示范工程正推进建设。

\subsection{核废物处理与安全保障}
核废物按放射性分为低中放废物（近地表处置）和高放废物（深地质处置）。我国建立了严格的核安全监管体系，三代核电技术实现“非能动安全”，大幅提升了核电安全性。

\section{可再生能源：能源转型的主力军}
\subsection{核心可再生能源的优缺点}
可再生能源是能源转型的核心，主要包括太阳能、风能、地热能：
- 太阳能：资源分布广、零污染，但受天气影响大、有间歇性；
- 风能：技术成熟、成本低，但同样存在间歇性，受地理条件限制；
- 地热能：稳定可靠、零排放，但资源分布不均、开发成本高。

\subsection{发展瓶颈与解决方案}
可再生能源的核心瓶颈是间歇性导致的并网消纳问题，解决方案包括：建设储能电站平抑波动、特高压输电实现远距离调度、源网荷储一体化提升系统灵活性。

\section{可燃冰：未来的潜力能源}
可燃冰是天然气水合物，能量密度高、储量巨大，但无法短期内取代石油：一是开采难度大，易引发甲烷泄漏、海底滑坡等环境风险；二是开采成本高，技术尚未成熟，仅处于试开采阶段。其定位是常规油气的补充，需稳步推进技术攻关与示范工程，防范环境风险。

\section{新型储能技术：新能源的“充电宝”}
储能是解决新能源间歇性的核心技术，相当于电网的“超大号充电宝”。主流技术包括锂电池、抽水蓄能、压缩空气储能等，未来将向高安全、长寿命、低成本、智能化方向发展。储能在新型电力系统中承担平抑波动、电网调峰、保障安全、支撑分布式能源等核心作用，是新能源大规模发展的关键支撑。

\section{新能源前沿技术与能源政策}
\subsection{新能源前沿技术}
课堂所学的前沿技术将推动能源转型：光催化制氢可生产“绿氢”，钍基熔盐堆提升核电安全性，钙钛矿太阳能电池有望替代传统硅电池，大数据技术可实现新能源智能调度。

\subsection{能源政策与市场改革}
碳排放交易市场通过市场化手段倒逼企业减排，“双碳”目标为能源转型指明方向。燃煤火电将向清洁高效、调峰辅助转型，能源市场改革培育多元主体，催生虚拟电厂、源网荷储一体化等新型商业模式，为能源转型提供制度保障。

\section{结论}
通过课程学习与报告梳理，我认识到能源转型是长期过程，化石能源将逐步转型，核能、可再生能源、储能将成为未来能源体系的核心。作为电气专业学生，我将继续学好专业知识，为我国双碳目标与新型电力系统建设贡献力量。

\section{参考文献}
\begin{enumerate}
    \item 《走进新能源》课程讲义
    \item 国家发展改革委, 国家能源局. “十四五”现代能源体系规划[R]. 2022.
    \item 国际能源署. 世界能源展望(2025)[R]. 2025.
\end{enumerate}

\end{document}